\slajdid{02-s016}
\begin{frame}[label=02-s016]
\gusttekst
\begin{columns}[c,onlytextwidth]
\begin{column}{.70\textwidth}
\alert{Drugi način} tretirajući funkciju \(\bar F\) na isti način kako kada smo posmatrili da funkcija F ima vrednosti jedan (funkcija \(\bar F\) će imati vrednost logičke jedinice tamo gde funkcija F ima vrednost logičke nule)
\[
\bar F=\bar C\bar B\bar A+\bar C\bar BA+\bar CB\bar A+C\bar B\bar A+CBA
\]
\end{column}
\begin{column}{.27\textwidth}
\begingroup
\setlength{\tabcolsep}{2pt}
\renewcommand{\arraystretch}{1.12}
\par\noindent
\begin{tabularx}{\linewidth}{*{5}{>{\centering\arraybackslash}X}}
\toprule C&B&A&F&\(\bar F\)\\\midrule
\textcolor{DEcrvena}{0}&\textcolor{DEcrvena}{0}&\textcolor{DEcrvena}{0}&\textcolor{DEcrvena}{0}&\textcolor{DEcrvena}{1}\\
\textcolor{DEcrvena}{0}&\textcolor{DEcrvena}{0}&\textcolor{DEcrvena}{1}&\textcolor{DEcrvena}{0}&\textcolor{DEcrvena}{1}\\
\textcolor{DEcrvena}{0}&\textcolor{DEcrvena}{1}&\textcolor{DEcrvena}{0}&\textcolor{DEcrvena}{0}&\textcolor{DEcrvena}{1}\\
0&1&1&1&0\\
\textcolor{DEcrvena}{1}&\textcolor{DEcrvena}{0}&\textcolor{DEcrvena}{0}&\textcolor{DEcrvena}{0}&\textcolor{DEcrvena}{1}\\
1&0&1&1&0\\
1&1&0&1&0\\
\textcolor{DEcrvena}{1}&\textcolor{DEcrvena}{1}&\textcolor{DEcrvena}{1}&\textcolor{DEcrvena}{0}&\textcolor{DEcrvena}{1}\\\bottomrule
\end{tabularx}\par
\endgroup

\end{column}
\end{columns}
\[
\begin{aligned}
\overline{(\bar F)}
&=\overline{(\bar C\bar B\bar A+\bar C\bar BA+\bar CB\bar A+C\bar B\bar A+CBA)}\\
&=\overline{\bar C\bar B\bar A}\cdot\overline{\bar C\bar BA}\cdot
\overline{\bar CB\bar A}\cdot\overline{C\bar B\bar A}\cdot\overline{CBA}
\end{aligned}
\]
\[
F=(C+B+A)(C+B+\bar A)(C+\bar B+A)(\bar C+B+A)(\bar C+\bar B+\bar A)
\]
\end{frame}
